\documentclass[10pt, letterpaper]{article}

%------------------------------------------------------------------------------
%  CMU MIIS Resume - Research-focused, URLs at bottom
%------------------------------------------------------------------------------

% Packages:
\usepackage[
    ignoreheadfoot,
    top=1.5 cm,
    bottom=1.5 cm,
    left=1.5 cm,
    right=1.5 cm,
    footskip=0.8 cm
]{geometry}
\usepackage{titlesec}
\usepackage{tabularx}
\usepackage{array}
\usepackage[dvipsnames]{xcolor}
\definecolor{primaryColor}{RGB}{0, 79, 144}
\usepackage{enumitem}
\usepackage{fontawesome5}
\usepackage{amsmath}
\usepackage[
    pdftitle={Zhaoxiang (Aaron) Feng's Resume},
    pdfauthor={Zhaoxiang Feng},
    colorlinks=true,
    linkcolor=primaryColor,
    urlcolor=primaryColor
]{hyperref}
\usepackage[pscoord]{eso-pic}
\usepackage{calc}
\usepackage{bookmark}
\usepackage{lastpage}
\usepackage{changepage}
\usepackage{paracol}
\usepackage{ifthen}
\usepackage{needspace}
\usepackage{iftex}

\ifPDFTeX
    \input{glyphtounicode}
    \pdfgentounicode=1
    \usepackage[utf8]{inputenc}
    \usepackage[T1]{fontenc}
    \usepackage{lmodern}
\fi

%------------------------------------------------------------------------------
%  Layout settings
%------------------------------------------------------------------------------
\AtBeginEnvironment{adjustwidth}{\partopsep0pt}
\pagestyle{empty}
\setcounter{secnumdepth}{0}
\setlength{\parindent}{0pt}
\setlength{\topskip}{0pt}
\setlength{\columnsep}{0cm}

\titleformat{\section}{\needspace{4\baselineskip}\bfseries\large}{}{0pt}{}[\vspace{1pt}\titlerule]
\titlespacing{\section}{-1pt}{0.20 cm}{0.15 cm}

\renewcommand\labelitemi{$\circ$}
\newenvironment{highlights}{
    \begin{itemize}[
        topsep=0.06 cm,
        parsep=0.06 cm,
        partopsep=0pt,
        itemsep=0pt,
        leftmargin=0.4 cm + 10pt
    ]
}{\end{itemize}}

\newenvironment{onecolentry}{
    \begin{adjustwidth}{0.2 cm + 0.00001 cm}{0.2 cm + 0.00001 cm}
}{\end{adjustwidth}}

\newenvironment{twocolentry}[2][]{
    \onecolentry
    \def\secondColumn{#2}
    \setcolumnwidth{\fill, 4.5 cm}
    \begin{paracol}{2}
}{
    \switchcolumn \raggedleft \secondColumn
    \end{paracol}
    \endonecolentry
}

\newenvironment{header}{
    \setlength{\topsep}{0pt}\par\kern\topsep\centering\linespread{1.5}
}{
    \par\kern\topsep
}

\newcommand{\AND}{\unskip\enspace\textbar\enspace}

\begin{document}

    %--------------------------------------------------------------------------
    %  Header
    %--------------------------------------------------------------------------
    \begin{header}
        \textbf{\fontsize{24 pt}{24 pt}\selectfont Zhaoxiang (Aaron) Feng}

        \vspace{0.25 cm}

        \normalsize
        \mbox{{\footnotesize\faMapMarker*}\hspace*{0.13cm}San Diego, CA}%
        \kern 0.25 cm\AND\kern 0.25 cm%
        \mbox{{\footnotesize\faEnvelope[regular]}\hspace*{0.13cm}\href{mailto:zhf004@ucsd.edu}{zhf004@ucsd.edu}}%
        \kern 0.25 cm\AND\kern 0.25 cm%
        \mbox{{\footnotesize\faPhone*}\hspace*{0.13cm}\href{tel:+13363377139}{(336)\,337-7139}}%
        \kern 0.25 cm\AND\kern 0.25 cm%
        \mbox{{\footnotesize\faGithub}\hspace*{0.13cm}\href{https://github.com/aaronzhfeng}{aaronzhfeng}}%
        \kern 0.25 cm\AND\kern 0.25 cm%
        \mbox{{\footnotesize\faGlobe}\hspace*{0.13cm}\href{https://aaronzhfeng.github.io}{aaronzhfeng.github.io}}%
    \end{header}

    \vspace{0.15 cm}

    %--------------------------------------------------------------------------
    %  Education
    %--------------------------------------------------------------------------
    \section{Education}

    \begin{twocolentry}{\textit{Sept.\,2022 -- Jun.\,2026 (expected)}}
        \textbf{University of California San Diego (UCSD)}

        \textit{B.S. in Data Science \& B.S. in Probability \& Statistics}
    \end{twocolentry}

    \vspace{0.06 cm}
    \begin{onecolentry}
        \begin{highlights}
            \item GPA: 3.83 (Major GPAs: 3.94 in Data Science, 3.93 in Statistics)
            \item Honors: Provost Honors (multiple quarters)
            \item Graduate Coursework: Advanced Time Series Analysis, Machine Learning, Applied Statistics, Privacy-sensitive Data Systems
        \end{highlights}
    \end{onecolentry}

    %--------------------------------------------------------------------------
    %  Publications
    %--------------------------------------------------------------------------
    \section{Publications}
    
    \begin{onecolentry}
        {[1]} Matthew Ho, Chen Si, \textbf{Zhaoxiang Feng}, Fangxu Yu, Yichi Yang, Zhijian Liu, Zhiting Hu, Lianhui Qin. ``ArcMemo: Abstract Reasoning Composition with Lifelong LLM Memory.'' \hfill [\href{https://arxiv.org/abs/2509.04439}{paper}] [\href{https://github.com/matt-seb-ho/arc\_memo}{code}]\\
        \textit{Under review at International Conference on Learning Representations (ICLR)}, 2026.\\
        \textit{\href{https://arcprize.org/blog/arc-prize-2025-results-analysis}{\textbf{Runner Up}}, ARC Prize 2025 Paper Awards (\$2,500) --- Top 8 of 90 paper submissions.}
    \end{onecolentry}

    \vspace{0.15 cm}

    \begin{onecolentry}
        {[2]} \textbf{Zhaoxiang Feng}, David Scott Lewis, Enrique Zueco. ``LabMemo: Concept-Level Memory for Autonomous Scientific Discovery.'' \hfill [\href{https://github.com/aaronzhfeng/lab\_memo/blob/main/paper/LabMemo\_Paper.pdf}{paper}] [\href{https://github.com/aaronzhfeng/lab\_memo}{code}]\\
        \textit{Accepted at IEEE IROS 2025 Workshop on Embodied AI and Robotics for Future Scientific Discovery (AIR4S)}, 2025.
    \end{onecolentry}

    %--------------------------------------------------------------------------
    %  Research Experience
    %--------------------------------------------------------------------------
    \section{Research Experience}

    % Q-Lab
    \begin{twocolentry}{Jan. 2025 -- Present}
        \textbf{Research Assistant | Q-Lab, UC San Diego}\\
        \textit{Advisors: Prof. Lianhui Qin and Matthew Ho}
    \end{twocolentry}
    \vspace{0.06 cm}
    \begin{onecolentry}
        \begin{highlights}
            \item Led design of program-synthesis-style memory ontology for ArcMemo, developing many-to-one puzzle-to-feature mappings and concept parameterizations for concept-level reasoning.
            \item Engineered reasoning-based retrieval mechanism achieving 7.5\% relative gain on ARC-AGI-1 (59.33\% official score); extended framework to AIME math problems with 9.3\% accuracy improvement.
            \item Built complete concept dataset generation pipeline transforming hand-written concepts into validated helper puzzles through multi-stage LLM-based generation and automated testing.
        \end{highlights}
    \end{onecolentry}
    \vspace{0.12 cm}

    % Wang Lab
    \begin{twocolentry}{Jun. 2025 -- Present}
        \textbf{Research Assistant | Wang Lab, UC San Diego}\\
        \textit{Advisors: Prof. Wei Wang and Young Su Ko}
    \end{twocolentry}
    \vspace{0.06 cm}
    \begin{onecolentry}
        \begin{highlights}
            \item Developing ProteinAgent: agentic framework orchestrating domain tools (AlphaFold, MD simulations) using concept-level memory to guide protein stability analysis protocols.
            \item Architected heterogeneous Mixture-of-Experts system for chemical reaction prediction, integrating graph-based, SMILES, 3D geometry-aware, and condition-aware expert models on USPTO datasets.
        \end{highlights}
    \end{onecolentry}
    \vspace{0.12 cm}

    % HDSI Capstone
    \begin{twocolentry}{Sep. 2025 -- Present}
        \textbf{Student Researcher | Halıcıo\u{g}lu Data Science Institute, UC San Diego}\\
        \textit{Advisor: Prof. Hao Zhang}
    \end{twocolentry}
    \vspace{0.06 cm}
    \begin{onecolentry}
        \begin{highlights}
            \item Senior Capstone: Building scalable pipelines for open-source LLM training, inference benchmarking, and alignment strategies.
        \end{highlights}
    \end{onecolentry}

    %--------------------------------------------------------------------------
    %  Teaching Experience
    %--------------------------------------------------------------------------
    \section{Teaching Experience}

    \begin{twocolentry}{\textit{Summer 2025}}
        \textbf{Tutor --- DSC 40A: Theoretical Foundations of Data Science I}\\
        \textit{University of California San Diego}
    \end{twocolentry}
    \vspace{0.06 cm}
    \begin{onecolentry}
        \begin{highlights}
            \item Conducted tutoring sessions covering probability and algorithmic thinking; led review sessions and supervised oral exams.
        \end{highlights}
    \end{onecolentry}

    %--------------------------------------------------------------------------
    %  Technical Skills
    %--------------------------------------------------------------------------
    \section{Technical Skills}

    \begin{onecolentry}
        \textbf{Languages \& Frameworks:} Python, PyTorch, PyTorch Geometric, Transformers, scikit-learn, R, SQL, Java\\[2pt]
        \textbf{LLM/ML:} OpenAI APIs, LLM fine-tuning (PPO, DPO), RAG systems, memory-augmented LLMs, reasoning-based retrieval, vision-language models, prompt engineering\\[2pt]
        \textbf{Infrastructure:} AWS, Docker, distributed training, Git/GitHub, Linux, vector databases
    \end{onecolentry}

\end{document}

